\documentclass[11pt]{article}

\usepackage[margin=1in]{geometry}
\usepackage{amsmath,amssymb,amsthm}
\usepackage{graphicx}
\usepackage[hidelinks]{hyperref}

\newtheorem{theorem}{Theorem}
\newtheorem{corollary}{Corollary}
\newtheorem{example}{Example}

\title{A Norming-Face Obstruction for \texorpdfstring{$L^p$}{Lp} State Spaces}
\author{Automatic Math Discovery Smoke Test}
\date{June 13, 2026}

\begin{document}
\maketitle

\section*{Source problem}

The source paper asks for an analogue, for $1<p<\infty$, of its
characterization of weakly compact state spaces in $L^1(\mu,X)$.

\begin{center}
\includegraphics[width=\textwidth]{figures/open_problem_crop.png}
\end{center}

\section*{Result}

This packet gives a necessary condition for weak compactness of the state
space of a constant function. It supplies a simple non-reflexive obstruction
which is not covered by the reflexive-domain partial answer.

Throughout this note the scalar field is real.

\begin{theorem}
Let $(\Omega,\Sigma,\mu)$ be a probability space, let $1<p<\infty$, and let
$q$ be conjugate to $p$. Let $X$ be a Banach space, let $x\in S_X$, and put
$f(\omega)=x$ for all $\omega\in\Omega$. Define
\[
S_f=\left\{g\in B_{L^q(\mu,X^*)}:
  \int_\Omega g(\omega)(x)\,d\mu(\omega)=1\right\}
\]
and
\[
F_x=\{x^*\in B_{X^*}:x^*(x)=1\}.
\]
If $S_f$ is weakly compact in $L^q(\mu,X^*)$, then $F_x$ is weakly compact in
$X^*$.
\end{theorem}

\begin{proof}
For each $x^*\in F_x$, let $Jx^*$ be the constant function with value
$x^*$. Since $\mu$ is a probability measure,
\[
\|Jx^*\|_{L^q(\mu,X^*)}=\|x^*\|\leq 1
\]
and
\[
\int_\Omega (Jx^*)(\omega)(x)\,d\mu(\omega)=x^*(x)=1.
\]
Thus $J(F_x)\subset S_f$.

Now define the bounded linear operator
\[
Q:L^q(\mu,X^*)\longrightarrow X^*,\qquad
Q(g)=\int_\Omega g(\omega)\,d\mu(\omega),
\]
where the integral is the Bochner integral. Holder's inequality gives
$\|Q(g)\|\leq \|g\|_{L^q(\mu,X^*)}$, and $Q(Jx^*)=x^*$ for every
$x^*\in F_x$.

If $S_f$ is weakly compact, then $J(F_x)$ is relatively weakly compact as a
subset of $S_f$. Since $Q$ is weak-to-weak continuous,
$Q(J(F_x))=F_x$ is relatively weakly compact in $X^*$. Finally, $F_x$ is
weakly closed: it is the intersection of the weakly closed unit ball
$B_{X^*}$ with the weakly closed affine hyperplane $\{x^*:x^*(x)=1\}$.
Therefore $F_x$ is weakly compact.
\end{proof}

\begin{corollary}
If $x\in S_X$ has a non-weakly-compact norming face $F_x$, then the state
space of the constant function $f(\omega)=x$ is not weakly compact in
$L^q(\mu,X^*)$.
\end{corollary}

\begin{example}
Take $X=\ell_1$ and $x=e_1$. Then $X^*=\ell_\infty$ and
\[
F_{e_1}=\{a=(a_n)\in B_{\ell_\infty}:a_1=1\}.
\]
This face is not weakly compact. Indeed, translating by $(-1,0,0,\ldots)$
identifies it with $\{0\}\times B_{\ell_\infty}$, which is isometric to the
closed unit ball of $\ell_\infty$. If $F_{e_1}$ were weakly compact, then
$B_{\ell_\infty}$ would be weakly compact, forcing $\ell_\infty$ to be
reflexive, a contradiction.

Hence, for any probability space and every $1<p<\infty$, the state space of
the constant $e_1$ function in $L^p(\mu,\ell_1)$ is not weakly compact.
\end{example}

\section*{Verification code}

The script \texttt{code/check\_constant\_face\_obstruction.py} verifies a
finite-truncation version of the constant-face inclusion: for several vectors
$a\in B_{\ell_\infty^n}$ with first coordinate $1$, the constant function has
$L^q$ norm $1$ and pairing $1$ with the constant $e_1$ function.

\section*{Scope}

This is a partial answer to Problem 3.18. It does not characterize all weakly
compact $L^p$ state spaces. It says that, even for constant functions, weak
compactness of the state space detects weak compactness of the corresponding
norming face in $X^*$.

\end{document}
